\section{Estado da Arte}

A fabricação de cerveja artesanal depende criticamente da conversão do amido em açúcares fermentescíveis, processo mediado por enzimas como a \textbf{\(\beta\)-amilase} e \textbf{\(\alpha\)-amilase} durante a etapa de mosturação. Tradicionalmente, o teste de iodo é utilizado para verificar a presença de amido residual no mosto, mas sua aplicação manual é subjetiva e pouco precisa, dependendo da interpretação visual do operador. Nesse contexto, a automatização desse teste pode melhorar significativamente a repetibilidade, precisão e eficiência do processo.

Este Estado da Arte analisa os estudos existentes sobre monitoramento da degradação do amido, identificando metodologias, limitações e oportunidades para a automatização do teste de iodo na fabricação de cerveja artesanal.

---

\subsection{Automação do Controle de Temperatura sem Monitoramento Bioquímico}
\cite{parizotto2017} desenvolveu um sistema de automação para controle de temperatura e tempo nas etapas de mosturação e fervura, utilizando um microcontrolador MSP-EXP432P401R, sensores PT-100 e controladores PI. O sistema demonstrou precisão no controle de temperatura, com tempo de assentamento de 10 minutos e erro nulo em regime permanente, além de um custo acessível (R\$ 2.626,30).

No entanto, o estudo não aborda o monitoramento da degradação do amido, etapa crítica para garantir a qualidade do mosto. A integração de um teste de iodo automatizado poderia complementar o sistema, permitindo ajustes em tempo real com base na atividade enzimática, otimizando a conversão do amido em açúcares fermentescíveis.

---

\subsection{Uso do Teste de Iodo para Monitoramento Manual da Degradação do Amido}
\cite{almeida2024} explorou o processo de fabricação de cerveja como ferramenta pedagógica para ensinar bioquímica, com foco na degradação enzimática do amido (amilose e amilopectina). O estudo utilizou testes manuais de iodo (2\%) para detectar amido residual, além de refratômetros para medir a concentração de açúcares e análises sensoriais para avaliar a conversão do amido.

Embora o teste de iodo tenha se mostrado eficaz para detectar amido residual, sua aplicação manual é subjetiva e pouco precisa para uso industrial ou em automação. Além disso, o estudo não aborda a automação do teste ou sua integração a sistemas de controle de processo, limitando seu potencial para aplicações práticas na fabricação de cerveja artesanal.

---

\subsection{Análise Comparativa e Lacunas Identificadas}
Os estudos analisados abordam aspectos complementares, mas não integrados no contexto da fabricação de cerveja artesanal. Enquanto \cite{parizotto2017} avança na automação do controle de temperatura, \cite{almeida2024} demonstra a importância do teste de iodo para monitorar a degradação do amido, porém de forma manual.

A Tabela \ref{tab:comparativo} resume as principais diferenças e lacunas entre os estudos:

\begin{table}[h]
\centering
\caption{Comparativo entre os estudos analisados.}
\label{tab:comparativo}
\begin{tabular}{|p{5cm}|p{5cm}|p{5cm}|}
\hline
\textbf{Aspecto} & \textbf{\cite{parizotto2017}} & \textbf{\cite{almeida2024}} \\ \hline
\textbf{Foco principal} & Automação do controle de temperatura & Ensino de bioquímica via teste de iodo manual \\ \hline
\textbf{Monitoramento do amido} & Não & Sim (manual, subjetivo) \\ \hline
\textbf{Automação} & Sim (temperatura e tempo) & Não \\ \hline
\textbf{Integração com controle de processo} & Parcial (sem feedback bioquímico) & Não \\ \hline
\textbf{Lacuna principal} & Falta de monitoramento da degradação do amido & Falta de automação do teste de iodo \\ \hline
\end{tabular}
\end{table}

A principal lacuna identificada é a ausência de um sistema que automatize o teste de iodo e o integre ao controle de processo. Isso representa uma oportunidade para inovação, permitindo ajustes em tempo real com base na atividade enzimática e melhorando a qualidade da cerveja artesanal.

---

\subsection{Justificativa para o Projeto Proposto}
Diante das lacunas identificadas, este projeto propõe o desenvolvimento de uma plataforma automatizada para o teste de iodo, integrada a sistemas de controle de temperatura. A solução será projetada para:

\begin{itemize}
    \item Eliminar a subjetividade: Utilizar sensores ópticos ou espectrofotômetros miniaturizados para detectar automaticamente a cor da reação (violeta = amido presente; amarela = amido ausente).
    \item Integrar ao controle de processo: Conectar os resultados do teste a um sistema de ajuste automático de temperatura ou tempo de mosturação.
    \item Ser acessível: Empregar componentes de baixo custo (ex.: Arduino, Raspberry Pi, sensores ópticos) para viabilizar a solução para pequenos produtores.
\end{itemize}

A automatização do teste de iodo permitirá:
\begin{itemize}
    \item Monitoramento em tempo real da degradação do amido.
    \item Ajustes automáticos no processo (ex.: prolongar a mosturação se houver amido residual).
    \item Melhoria na qualidade da cerveja, com maior precisão e repetibilidade.
\end{itemize}

---

\subsection{Tendências Recentes: Aprendizado de Máquina e Inteligência Artificial}
Estudos recentes têm incorporado técnicas de aprendizado de máquina e inteligência artificial para análise e modelagem de sistemas complexos \cite{zhang2025}. Essas tecnologias podem ser aplicadas para:
\begin{itemize}
    \item Otimizar parâmetros de processo (ex.: temperatura ideal para ativação enzimática).
    \item Prever a qualidade da cerveja com base em dados históricos de degradação do amido e temperatura.
    \item Automatizar a interpretação dos resultados do teste de iodo, utilizando algoritmos de visão computacional para analisar a cor da reação.
\end{itemize}

A integração de IA e aprendizado de máquina ao sistema proposto pode expandir significativamente suas capacidades, tornando-o mais robusto e adaptável a diferentes condições de produção.

---

\subsection{Conclusão}
Este Estado da Arte revelou que, embora existam avanços no controle de temperatura \cite{parizotto2017} e no uso do teste de iodo para ensino \cite{almeida2024}, nenhum estudo automatiza esse teste ou o integra a sistemas de monitoramento. O projeto proposto preenche essa lacuna, desenvolvendo uma solução que combina automatização do teste de iodo com controle de processo, contribuindo para a inovação na fabricação de cerveja artesanal.

A incorporação de técnicas de IA e aprendizado de máquina \cite{zhang2025} pode ainda potencializar a plataforma, permitindo análises preditivas e otimização contínua do processo de mosturação.

---
